% !TeX spellcheck = hu_HU
% !TeX encoding = UTF-8
% !TeX program = xelatex
%----------------------------------------------------------------------------
\section{Primitív feladat ütemező}\label{sec:uses:make}
%----------------------------------------------------------------------------

A make rendszer egy elterjedt eszköz fordítási feladatok ütemezésére egy C vagy C++ projekt fordítása során.
Ehhez a make rendelkezik egy saját deklaratív \gls{DSL}-el, amelyiken belül meg lehet adni a feladatok közötti függőségeket.
Ezt szokásos a projekthez tartozó {\tt Makefile}-ba írni, esetleg generáltatni, ha komplexebb lépésekre van szüksége a projekt fordítása előtt, mint például platform képességeinek ellenőrzése.

Az előbb leírt felhasználási módra lett kitalálva a make.
Azonban a tényleges lépéseket, amit egy-egy feladat megvalósításához kell végrehajtani a rendszerhéj (system shell) értelmezője számára továbbítja, így sok minden egyéb feladatra lehet felhasználni, ami nem projektek fordítása, de olyan alfeladatokból áll, amelyek között függőségi sorrendet kell felállítani.

Ennek megfelelően, most bemutatok egy olyan primitív make jellegű \gls{DSL}-t, amelyik lehetővé teszi azt, hogy a egyes részfeladatok közötti függőségeket definiálva lehessen megvalósítani egy nagyobb feladatot.
A \gls{DSL}-t megvalósító C4 kód megtalálható \aref{sec:app:make}.~függelékben, a használatra adott példa pedig megtekinthető \aref{lst:ex-make}.~listázásban.
A példakódban az egyszerűség kedvéért csak egyszerű kiírások a feladatok, de ezek lehetnének komplexebb kódrészletek.

\begin{lstlisting}[gobble=4,numbers=left,caption={C4 make jellegű \gls{DSL} használati példa},label=lst:ex-make]
    let pipeline make
        rule "lint": {
            println "linting"
        }
        rule "compile": {
            println "compile"
        }
        rule "test" <: "lint", "compile": {
            println "test"
        }
        rule "package" <: "test": {
            println "package"
        }
    done_make

    run pipeline "package"
\end{lstlisting}

Hasonlóan az előző fejezetben tárgyalt állapotgép leírásához, a {\tt run} parancs az, amelyik a futtatáshoz tartozó szemantikát írja le a megadható feladatok számára.
Fontos megemlíteni, hogy ez a demó jellegű \gls{DSL} csak a feladatok ütemezését valósítja meg, az egyes fájl írás alapú késznek tekintettséget nem.
Bár ezt alapvetően nem lenne kifejezetten nehéz megvalósítani, a C4 megvalósításom jelenleg még nem rendelkezik az operációs rendszerrel ilyen módú interakció végrehajtására szükséges elemekkel.





